\section{Composition With Sibling LPs}
\label{sec:composition}

\subsection{LP-217 cert profile modes and the P3Q kinds}

The LP-217 cert-profile-mode ladder is the operator-facing
naming for cross-family PQ posture. Each mode declares a
required leg set; the P3Q kinds supply the on-chain machinery
that makes each mode's leg set callable from a rollup contract:

\begin{table}[h]
\centering
\small
\begin{tabular}{@{}lp{4.5cm}rl@{}}
\toprule
\textbf{Mode} & \textbf{Required legs (P3Q calls)} & \textbf{Per-batch gas} & \textbf{Use case} \\
\midrule
\texttt{PQ-off} & BLS only (no P3Q) & $\sim$\SI{5000}{} & non-financial; gaming \\
\texttt{PQ-fast} & BLS + P3Q($0x01$) & $\sim$\SI{55000}{} & payments, DEX, gaming \\
\texttt{PQ-strict} & BLS + P3Q($0x01$) + P3Q($0x02$) & $\sim$\SI{5055000}{} & custody, treasury, securities \\
\textbf{\texttt{PQ-heavy}} & BLS + P3Q($0x01$) + P3Q($0x02$) + P3Q($0x03$) & $\sim$\SI{5085000}{} & sovereign anchors, bridges \\
\bottomrule
\end{tabular}
\caption{LP-217 modes mapped to P3Q kind calls. The four cert
profile modes correspond to subsets of the P3Q kind set; PQ-heavy
calls all three. The $\sim$\SI{5}{M}-gas cost cliff between
PQ-fast and PQ-strict is the price of intra-lattice parameter-diversity
insurance via the Corona kind.}
\label{tab:cert-modes-p3q}
\end{table}

The composition is \textbf{additive}: a PQ-heavy rollup contract
calls the BLS precompile and three \texttt{P3Q.verify(kind, ...)}
calls at slot \texttt{0x012205} in sequence inside its
\texttt{validateBatch} function. If any of the four returns
\texttt{false}, the batch is rejected.

\subsubsection{Sample rollup-contract composition (Solidity)}

\begin{lstlisting}[language={}]
function validateBatch(
    bytes32 prevRoot,
    bytes32 newRoot,
    bytes32 batchHash,
    bytes calldata blsSig,
    bytes calldata pulsarSig,
    bytes calldata coronaSig,    // omit for PQ-strict
    bytes calldata magnetarSig,  // omit for PQ-strict/fast
    bytes calldata groupPk
) external returns (bool) {
    bytes32 h = sha256(abi.encodePacked(prevRoot, newRoot, batchHash));

    // Classical leg
    if (!BLS.verify(blsSig, h, blsGroupPk)) return false;

    // P3Q Pulsar leg (PQ-fast and above)
    if (!P3Q.verify(0x01, pulsarSig, h, groupPk)) return false;

    if (mode >= PQ_STRICT) {
        // P3Q Corona leg (intra-lattice parameter/impl diversity)
        if (!P3Q.verify(0x02, coronaSig, h, groupPk)) return false;
    }

    if (mode == PQ_HEAVY) {
        // P3Q Magnetar leg (hash-based family independence)
        if (!P3Q.verify(0x03, magnetarSig, h, groupPk)) return false;
    }

    return true;
}
\end{lstlisting}

The same precompile address (\texttt{0x012205}) is called
three times with three different kind bytes. The rollup
contract's mode field is set at rollup-genesis time via
\texttt{CertModeFloor} (LP-218); the contract enforces it on
every batch.

\subsection{LP-218 rollup-batch pattern}

LP-218 specifies the rollup-batch envelope and the
sequencer$\to$Z-Chain submission pattern; LP-219 specifies the
per-kind verifier mechanics. The two LPs share the same slot
(\texttt{0x012205}), the same ABI shape
(\texttt{P3Q.verify(kind, sig, h, pk)}), and the same
\texttt{abiFalse}-on-cryptographic-failure contract.

LP-218 \S~``Precompile wire format'' documents the structural
format; the per-kind LPs LP-223, LP-220, and LP-222
document the per-kind specifics on top of that uniform skeleton.
The intent is clean cross-reference rather than duplication: a
reader who wants the rollup-batch pattern reads LP-218; a reader
who wants per-kind verifier mechanics reads LP-219; a reader who
wants both reads both.

\subsection{LP-300 schema registry}

The P3Q family lives at one slot, but the wire envelopes -- the
sequencer$\to$Z-Chain \texttt{RollupBatchTx} payload, the
inner-batch Pulsar/Corona/Magnetar signature blobs -- are
schema-versioned per LP-300. The registry entry for the P3Q
family is:

\begin{verbatim}
slot:    0x012205
kinds:   { 0x01: pulsar, 0x02: corona, 0x03: magnetar }
domain:  lux-evm-precompile-p3q-v1   (FIPS 204 ctx)
abi:     P3Q.verify(uint8, bytes, bytes32, bytes) -> bool
\end{verbatim}

Schema bumps within a kind (e.g.\ a future Corona parameter set
revision) advance the mode byte without changing the slot or
the kind byte. Schema bumps across the family (a fourth kind,
say) advance the kind byte without changing the slot.

\subsection{Cross-references}

\begin{table}[h]
\centering
\small
\begin{tabular}{@{}lp{9cm}@{}}
\toprule
\textbf{LP} & \textbf{Relationship} \\
\midrule
LP-073 & Pulsar threshold signature -- P3Q kind $0x01$ verifies \\
LP-181 & Magnetar threshold SLH-DSA -- P3Q kind $0x03$ verifies \\
LP-182 & QuasarCert -- the cert legs whose verify kinds LP-219 supplies \\
LP-200 & ZAP Stack umbrella -- the ``one and only one way'' principle the unified slot honors \\
LP-202 & Cert tier degradation -- the per-kind verify path participates in the tier ladder \\
LP-203 & GPU-native verify -- batched Corona / Magnetar kernels (projected) \\
LP-217 & Cert profile modes -- the operator-facing names that determine which P3Q kinds a rollup contract calls \\
LP-218 & ZAP-native PQ rollups -- specifies the unified slot and kind $0x01$ Pulsar; LP-219 specifies kinds $0x02$ Corona and $0x03$ Magnetar on the same slot \\
LP-0120 & Quasar mainnet defaults -- owns slots \texttt{0x012204},
  \texttt{0x012206}, \texttt{0x012207} for raw-primitive verifiers
  (the non-rollup-batch use case); LP-218/LP-219 P3Q family
  lives exclusively at \texttt{0x012205} \\
LP-300 & Schema registry -- entry for the P3Q family kind set \\
FIPS 204 & ML-DSA standard -- Pulsar-kind verifier conformant \\
FIPS 205 & SLH-DSA standard -- Magnetar-kind verifier conformant \\
Ringtail eprint 2024/1113 & 2-round lattice threshold reference -- Corona-kind verifier conformant \\
\bottomrule
\end{tabular}
\caption{LP-219 composition references.}
\label{tab:composition-refs}
\end{table}

\subsection{Backwards compatibility}

None. The pre-Quasar Edition Lux network had no P3Q precompile
family and is a separate network out of scope. The genesis state
of the new final Lux network (activation
\texttt{2025-12-25T16:20:00-08:00}) ships the unified-slot
kind-byte-dispatch ABI as the only ABI; no shim exists for the
sibling-slots draft.

\subsection{Future work}

\begin{itemize}[leftmargin=2em,nosep]
\item \textbf{Corona/Magnetar verifier wiring.} The reference
  \texttt{p3q.Run} returns \texttt{abiFalse} for kinds $0x02$
  and $0x03$ today. The follow-up is to wire
  \texttt{luxfi/corona/sign.Verify} and
  \texttt{luxfi/magnetar/verify} into the dispatch table. The
  ABI and the wire format are stable across this transition;
  only verify behavior changes from \texttt{false} to the
  cryptographic outcome.
\item \textbf{GPU dispatch boundary closure.} Install
  \texttt{lux-lattice.pc} on Spark and expose
  \texttt{lux\_slhdsa\_*} / \texttt{lux\_corona\_*} entry
  points in \texttt{lux/accel/c\_api.h}. Once both close, the
  Blackwell verify projections (LP-220
  and~LP-222) become measurements.
\item \textbf{Fourth PQ family.} If a code-based signature
  scheme (HQC-Sig variants, McEliece-Sig) becomes
  operationally relevant, it lands as one new
  \texttt{P3QVerifier} impl and one new \texttt{kind} byte
  (provisionally $0x04$) in the dispatch table. No new
  precompile slot allocation, no new Solidity function shape.
\item \textbf{Blackwell-discounted Corona gas.} Operators with
  all-Blackwell validator sets may petition for a discounted
  Corona gas charge reflecting GPU-class verify cost rather
  than CPU-class; LP-219 does not codify this. A future
  amendment would tie the discount to a validator-hardware
  attestation surface.
\end{itemize}
